\section{Execution-Model Contract}

\subsection{Why execution is part of the estimand}

For a passive maker, the fill set is generated jointly by market events and the
simulator's rules. Changing arrival timing, cancellation timing, queue
advancement, or post-only handling can change which trades fill an order. It can
therefore change P\&L, markouts, conditional signal tests, and apparent latency
sensitivity. Execution assumptions are not an implementation footnote; they
define the sample on which the strategy is evaluated.

L2 data reports aggregate quantity at a price, not individual order identities.
No replay using this dataset observes true queue rank. The implementation is a
documented execution approximation, not an exchange twin.

\subsection{Historical and current semantics}

\begin{table}[H]
\centering
\caption{Execution-model boundary.}
\label{tab:model-comparison}
\begin{tabularx}{\textwidth}{@{}L{0.25\textwidth}YY@{}}
\toprule
Feature & \LegacyModel & \CurrentModel \\
\midrule
New-order arrival & Eligible on the next depth update after modeled arrival & Scheduled at the exact arrival millisecond on the model clock \\
Cancellation & Effective immediately when requested & Scheduled separately; order stays fillable while cancel is in flight \\
Cancel/fill race & Absent & Explicit; a fill can beat a late cancel \\
Equal-time private action & Implicit through depth activation & All recorded market data at the millisecond precedes private arrivals \\
Own orders at one price & Incomplete separation of public and own queue & Non-overlapping own FIFO ranks with trade-volume conservation \\
Gap state & Strict replay policy added, but legacy artifacts retain old private semantics & Whole timestamp group invalidated fail-closed until snapshot resync \\
Snapshot timing & Returned state applied at a pre-fetch local tag & Local request-time resync barrier; recovery gated by post-response receipt order and released on depth \code{E}; new captures store request and response times \\
Artifact root & \path{results/panels/btcusdt_l2_panel_v2} & \path{results/panels/btcusdt_l2_panel_v3_event_driven} \\
Research status & Frozen historical evidence & Mechanics accepted; economic rerun pending \\
\bottomrule
\end{tabularx}
\end{table}

\subsection{Clock and equal-timestamp policy}

The replay clock has two classes of event. Its recorded-data axis combines local
request time for fail-closed resync barriers, depth \code{E} for ordinary diffs
and recovered-state release, and trade \code{T} for trades:

\begin{enumerate}
  \item recorded market data, ordered depth before trade for unresolved
        millisecond ties and preserving source order within each stream; and
  \item simulated private arrivals for new orders and cancellations, preserving
        stable insertion order.
\end{enumerate}

Private actions strictly before a market-data timestamp are processed first.
At an identical millisecond, the complete recorded market-data group is
processed before private arrivals. A new order arriving at the same millisecond
as a trade is therefore not eligible for that trade; a cancel arriving at the
same millisecond loses the race. This is a conservative policy for unresolved
timestamps, not a statement about the venue's internal matching sequence.

\begin{figure}[H]
\centering
\begin{tikzpicture}[
  event/.style={draw=Navy,rounded corners=2pt,fill=PaleBlue,align=center,
                minimum width=2.55cm,minimum height=0.85cm,font=\small},
  private/.style={draw=Orange,rounded corners=2pt,fill=PaleOrange,align=center,
                  minimum width=2.55cm,minimum height=0.85cm,font=\small},
  arrow/.style={-{Latex[length=2mm]},thick,draw=Slate}
]
  \draw[thick,Slate] (0,0) -- (12.8,0);
  \node[below] at (0,0) {$t-1$};
  \node[below] at (6.4,0) {$t$};
  \node[below] at (12.8,0) {$t+1$};
  \node[private] (before) at (2.0,1.0) {Private arrival\\strictly before $t$};
  \node[event] (depth) at (5.0,1.0) {Depth events\\at $t$};
  \node[event] (trade) at (7.8,1.0) {Trade events\\at $t$};
  \node[private] (equal) at (10.7,1.0) {Private arrivals\\at $t$};
  \draw[arrow] (before) -- (depth);
  \draw[arrow] (depth) -- (trade);
  \draw[arrow] (trade) -- (equal);
\end{tikzpicture}
\caption{\CurrentModel\ timestamp policy. The horizontal placement inside
the $t$ group shows processing precedence, not sub-millisecond observation.}
\label{fig:timestamp-policy}
\end{figure}

Depth-decrease attribution is deferred until every trade at the same
millisecond has been observed. Otherwise the depth-before-trade rule could
mistake traded quantity for a cancellation and grant queue improvement twice.

\subsection{Order and cancellation lifecycle}

Submitting an \code{OrderRequest} creates a pending order and schedules its
arrival at
\begin{equation}
  t_{\mathrm{arrive}} = t_{\mathrm{submit}} + L_e + J_e,
\end{equation}
where $L_e$ is base entry latency and $J_e$ is a seeded uniform integer jitter
bounded so that delay cannot be negative. Cancellation uses a separate latency
and independent random stream:
\begin{equation}
  t_{\mathrm{cancel}} = \max\left(t_{\mathrm{request}} + L_c + J_c,
  t_{\mathrm{arrive}}\right).
\end{equation}

While entry is pending, its quantity is counted as potential exposure but it is
not eligible for market fills. Once active, the order remains fillable until
the cancellation reaches the simulated exchange. Lifecycle events distinguish
request from effect:
\code{cancel_requested}, \code{cancelled}, and \code{cancel_too_late}. A
post-only limit that would cross on arrival is cancelled with an explicit
reason rather than silently reclassified as a passive fill.

\subsection{Queue model and own-order FIFO}

When a resting order arrives, its modeled quantity ahead is decomposed as
\begin{equation}
  Q^{\mathrm{ahead}}_i = Q^{\mathrm{external}}_i + Q^{\mathrm{own}}_i.
\end{equation}
The external component starts from displayed L2 quantity, subject to a FIFO
floor imposed by earlier live own orders. The own component is the remaining
quantity of earlier simulated orders at the same side and price. A public depth
decrease not explained by recorded trades improves only the external component:
\begin{equation}
  Q^{\mathrm{external}}_{i,t+1}
  = Q^{\mathrm{external}}_{i,t}\left(1-c\rho_t\right),
  \qquad c\in[0,1],
\end{equation}
where $c$ is queue-cancellation credit and $\rho_t$ is the inferred fraction of
the displayed queue removed without matching trade volume. The two historical
endpoints are $c=0$ and $c=1$; neither is claimed to be calibrated ground truth.

Recorded trade quantity advances non-overlapping FIFO ranks. The simulator
plans and validates the complete transition before mutating any order and
asserts
\begin{equation}
  \sum_i q^{\mathrm{maker\ fill}}_{i,t}
  \le q^{\mathrm{recorded\ trade}}_t.
\end{equation}
Cancelling an earlier own order releases its unfilled quantity from later own
queues. Decimal comparisons use only a scale-aware $10^{-24}$ relative
tolerance for division residue; the tolerance cannot create fill budget.

\subsection{Risk envelope}

Position is not the only relevant exposure during cancel/replace. If multiple
same-side entries are pending or cancels are in flight, all may still fill. The
strategy therefore checks worst-case working exposure:
\begin{equation}
  \left|I_t + \sum_{i\in\mathcal{F}_t^{\mathrm{same}}}
  s_i q^{\mathrm{remaining}}_i + s_{\mathrm{new}}q_{\mathrm{new}}\right|
  \le I_{\max},
\end{equation}
where $\mathcal{F}_t^{\mathrm{same}}$ includes active orders, pending entries
that may become fillable, and active orders with cancels in flight. An invariant
breach aborts replay; it is not recorded as a harmless temporary overshoot.

\subsection{Gaps and replay termination}

A known sequence gap invalidates the entire timestamp group before any event in
that group can fill an order or invoke strategy logic. Local orders and pending
actions become invalid because their queue state is no longer trustworthy.
Replay resumes only at a valid snapshot. This is fail-closed research
bookkeeping, not a claim that the exchange cancelled those orders instantly.

Private events scheduled after the last recorded market event are not executed
against an assumed future book. Remaining live orders are terminalized as
expired, and pending action counts are reported separately.

\subsection{Known execution limits}

\begin{itemize}
  \item Aggregate L2 cannot reveal true market-by-order rank, hidden liquidity,
        or where cancellations occurred relative to the simulated order.
  \item The modeled clock mixes recorded local request time for resync barriers,
        depth event time \code{E} for release and ordinary diffs, trade matching
        time \code{T}, and local receipt order used to gate snapshots.
        Those fields are not guaranteed to be perfectly comparable, and
        millisecond resolution cannot identify true within-millisecond matching
        order. The post-response receipt gate removes request-start placement,
        but its selected boundary remains on exchange \code{E}; it does not
        calibrate client observation or feed latency.
  \item Market-data latency is not modeled separately from private-message
        latency.
  \item Counterfactual taker execution shares a shadow liquidity ledger only
        until the next observed depth update; taker-heavy research needs market
        impact and stronger book reconciliation.
  \item No execution model here has been calibrated against live order
        acknowledgements and fills from an exchange account.
\end{itemize}
